% ============================================================
% Abstract -- trilingual, single page, small font (Juanola style)
% ============================================================
%\chapter*{Abstract}
\addcontentsline{toc}{chapter}{Abstract}
\thispagestyle{plain}

\begin{small}

\noindent\textbf{Abstract}

\smallskip
\noindent This thesis presents the design, implementation, and pilot validation of a
practical cybersecurity teaching package for the \textit{Introduction to
Cybersecurity} course at TecnoCampus. The package consists of four progressive
laboratory exercises built on EVE-NG, a network emulation platform deployed on
self-hosted infrastructure. Each laboratory provides a realistic multi-system
network topology, a structured student guide, and a teacher resolution guide with
assessment rubric. The four labs cover network reconnaissance and enumeration,
web application exploitation (SQL injection, cross-site scripting, authentication
bypass, and remote code execution), incident response and log forensics, and
cryptography and steganography. The full environment has been validated through a
pilot study with students enrolled in the course, confirming the reachability of
all learning objectives and the pedagogical adequacy of the materials.
\begin{comment}


\medskip
\noindent\textbf{Keywords:} cybersecurity education, penetration testing, EVE-NG,
network emulation, laboratory design, CTF, web security, incident response,
cryptography.
\end{comment}


\bigskip
\noindent\textbf{Resum}

\smallskip
\noindent Aquesta tesi presenta el disseny, la implementació i la validació pilot d'un paquet
didàctic pràctic de ciberseguretat per a l'assignatura \textit{Introducció a la
Ciberseguretat} del TecnoCampus. El paquet consisteix en quatre laboratoris
progressius construïts sobre EVE-NG, una plataforma d'emulació de xarxes desplegada
en infraestructura pròpia. Cada laboratori proporciona una topologia de xarxa
multi-sistema realista, una guia estructurada per a l'estudiant i una guia de
resolució per al professor amb rúbrica d'avaluació. Els quatre laboratoris cobreixen
reconeixement i enumeració de xarxes, explotació d'aplicacions web (injecció SQL,
XSS, evasió d'autenticació i execució remota de codi), resposta a incidents i
anàlisi forense de registres, i criptografia i esteganografia. L'entorn ha estat
validat mitjançant un estudi pilot amb estudiants matriculats a l'assignatura,
confirmant l'assolibilitat dels objectius d'aprenentatge i l'adequació pedagògica.
\begin{comment}
\medskip
\noindent\textbf{Paraules clau:} educació en ciberseguretat, proves d'intrusió,
EVE-NG, emulació de xarxes, disseny de laboratoris, CTF, seguretat web,
resposta a incidents, criptografia.
\end{comment}


\bigskip
\noindent\textbf{Resumen}

\smallskip
\noindent Esta tesis presenta el diseño, la implementación y la validación piloto de un
paquete didáctico práctico de ciberseguridad para la asignatura \textit{Introducción
a la Ciberseguridad} del TecnoCampus. El paquete consiste en cuatro laboratorios
progresivos construidos sobre EVE-NG, una plataforma de emulación de redes
desplegada en infraestructura propia. Cada laboratorio proporciona una topología
de red multi-sistema realista, una guía estructurada para el estudiante y una
guía de resolución para el profesor con rúbrica de evaluación. Los cuatro
laboratorios cubren reconocimiento y enumeración de redes, explotación de
aplicaciones web (inyección SQL, XSS, evasión de autenticación y ejecución remota
de código), respuesta a incidentes y análisis forense de registros, y criptografía
y esteganografía. El entorno ha sido validado mediante un estudio piloto con
estudiantes matriculados en la asignatura, confirmando la alcanzabilidad de los
objetivos de aprendizaje y la adecuación pedagógica de los materiales.
\begin{comment}
\medskip
\noindent\textbf{Palabras clave:} educación en ciberseguridad, pruebas de
intrusión, EVE-NG, emulación de redes, diseño de laboratorios, CTF, seguridad
web, respuesta a incidentes, criptografía.
\end{comment}
\end{small}
\cleardoublepage
